一行代码写着 \texttt{loss.backward()}，它算的是 \(\nabla\E[\ell]\) 还是 \(\E[\nabla\ell]\)？

框架不会告诉你。它把梯度算子塞进期望里，然后照常返回一个张量。多数时候这样做是对的，因为该成立的条件确实成立；但条件不成立时它也不报错——你会得到一个看上去正常、实际有偏的梯度。重参数化技巧成立的边界、\(\Uniform(0,\theta)\) 这种支撑随参数移动的情形、以及蒙特卡洛估计里偶发的跳数，全都落在这条缝上。

这一部分处理的就是这类问题：\textbf{什么时候可以交换两个操作的顺序。}极限与积分、导数与期望、求和与求和——这些交换在前面各部分被反复使用，而它们成立的条件从未被完整交代过。

追到底，所有这些条件都在防同一件事：\emph{质量逃逸}。函数列逐点收敛到零，积分却恒为 1，因为面积跑到了越来越窄的尖峰里；期望的极限不等于极限的期望，因为概率质量溜到了尾部。控制收敛定理里那个可积包络，紧致性里那个有限覆盖，一致可积性里那个一致的界，做的都是同一件事——把质量按住，不让它跑掉。

\begin{center}
\begin{tikzpicture}[x=1cm,y=1cm,font=\small,>=stealth]
  \draw[black!45,fill=blue!9] (0,1.4) rectangle (3.6,2.9);
  \node[align=center] at (1.8,2.15) {想交换\\两个操作};
  \draw[black!45,fill=red!7] (5.1,3.5) rectangle (13.3,4.4);
  \node[align=left,anchor=west] at (5.3,3.95) {风险：质量逃逸——尖峰、尾部、支撑漂移};
  \draw[black!45,fill=green!8] (5.1,2.35) rectangle (13.3,3.25);
  \node[align=left,anchor=west] at (5.3,2.8) {按住它的条件：可积包络、一致有界、紧致};
  \draw[black!45,fill=green!8] (5.1,1.2) rectangle (13.3,2.1);
  \node[align=left,anchor=west] at (5.3,1.65) {记账的语言：测度、可测函数、\(L^p\)};
  \draw[black!45,fill=green!8] (5.1,0.05) rectangle (13.3,0.95);
  \node[align=left,anchor=west] at (5.3,0.5) {几何的语言：流形、度量、投影};
  \draw[->,black!55] (3.65,2.3) -- (5.05,3.95);
  \draw[->,black!55] (3.65,2.2) -- (5.05,2.8);
  \draw[->,black!55] (3.65,2.05) -- (5.05,1.65);
  \draw[->,black!55] (3.65,1.9) -- (5.05,0.5);
\end{tikzpicture}

\emph{这一部分不生产新的计算方法，它生产的是许可证：说明前面用过的那些操作，各自在什么条件下合法。}
\end{center}

\subsection{五条贯穿始终的判断}

\textbf{稠密不等于没有洞。}有理数在数轴上处处稠密，却漏掉了 \(\sqrt2\)。完备性把这些洞补上，于是``逼近序列的极限确实存在且仍在空间里''成为可以依赖的事实——一切迭代算法的收敛论证、一切存在性定理，都建立在这一条上。Cauchy 判据的好处是不必先知道极限：看序列自己后面的项挤不挤在一起就够了。

\textbf{逐点成立与一致成立差一个量词顺序，而这个差别决定一切。}点态收敛允许每个点有自己的 \(N\)，一致收敛要求一个 \(N\) 对所有点通用。前者不保持连续性、不能与积分交换；后者可以。同一个区分在紧致性那里换了个面孔——紧致的作用正是把无限个局部结论合并成一个整体结论，让你能从无限多个正数里取到一个正的最小值。

\textbf{\(\sigma\) 代数是``哪些问题可以被回答''的写法。}可测集就是能谈概率的集合；日志的聚合粒度决定了哪些问题日后还问得出来。而条件期望正是对信息 \(\sigma\) 代数的投影——这一句把回归、滤波、鞅统一起来，也解释了为什么残差与特征正交是投影的定义性质而不是拟合好坏的证据。

\textbf{密度是相对于某个基准测度才有意义的概念。}同一个分布相对 Lebesgue 测度有密度，相对计数测度就没有；换个单位，密度值整体平移几个数量级。Radon–Nikodym 导数把``相对于谁''说清楚，重要性采样的权重、似然比、KL 散度都是它的实例。

\textbf{几何结构是外加的，不是对象自带的。}同一个流形可以配不同的黎曼度量，得到完全不同的距离与最速下降方向。Fisher 度量只是统计流形上的一种装法——但它是唯一一种在充分统计量变换下不变的装法，这解释了自然梯度为什么值得算。换尺子会换掉所有下游结论，所以尺子必须被显式声明。

\subsection{十一章怎么安排}

这一部分不必按顺序通读。多数读者是带着一个具体问题来的，下面按问题给出入口。

想知道\emph{极限与积分何时可交换}，读\hyperref[ch:100]{``什么时候可以交换极限与运算''}——它是本部分最直接有用的一章，重参数化梯度、蒙特卡洛估计的合法性都在这里。前置只需要\hyperref[ch:098]{``从数列到函数列的收敛''}里点态与一致的区分。

想给概率论回填地基，读\hyperref[ch:101]{``测度、积分与条件期望''}。它被前面多个部分引用，其中条件期望作为投影那一篇价值最高。\hyperref[ch:102]{``乘积测度与 Fubini''}紧随其后，处理独立性与累次积分换序。

做优化或深度学习的读者，\hyperref[ch:105]{``信息几何：参数空间中的分布几何''}回报最直接：它解释了参数空间的欧氏距离为何没有意义、自然梯度在算什么、以及信任域方法里 KL 约束的几何含义。前置是\hyperref[ch:104]{``流形初步''}的前三节。

\hyperref[ch:097]{``实数为什么需要完备''}、\hyperref[ch:099]{``度量、紧致与完备''}、\hyperref[ch:103]{``Lp 空间与函数几何''}构成分析的骨架，按需回补。\hyperref[ch:107]{``曲线与曲面几何''}相对独立，其中 Gauss 绝妙定理说明了一个看似关于嵌入方式的量其实纯粹内蕴。

\hyperref[ch:106]{``本质结构：对称性、低维性、信息几何与因果''}是全书的交汇章，建议放到最后读——它把对称性与等变、低维结构、信息几何、因果图这四条线各自的依据摆在一起，并明确区分哪些是定理、哪些是经验观察、哪些只是启发。

读这一部分时值得记住：它给出的不是新方法，而是前面那些方法的适用条件。条件写在纸上时看着像形式主义，直到某次结果不对、而你需要知道该去查哪一条。
